\begin{abstract}
  \noindent Tässä tutkielmassa arvioidaan kielimallipohjaisen Text-to-SQL-järjestelmän
  kykyä muodostaa käyttäjän tiedontarvetta vastaavia SQL-kyselyitä sekä
  tulkinnan johdonmukaisuutta kysymysten toistoissa ja sanamuodon
  muutoksissa.

  Teoreettisessa osassa jäsennetään se matemaattinen ja arkkitehtuurinen jatkumo, jolla diskreeteistä tekstiesityksistä, optimointimenetelmistä ja itsehuomiomekanismista rakentuu symbolista koodia tuottava autoregressiivinen todennäköisyysmalli. Teoria tarjoaa käsitteellisen perustan kielimallin toiminnalle ja havaitun vaihtelun ymmärtämiselle.

  Empiirinen osa on eksploratiivinen pilottitutkimus Snowflaken
  Cortex Analyst -järjestel\-mästä, joka hyödyntää semanttista näkymää
  puhelinkeskusdataa sisältävässä tietokantaympäristössä.

  Tutkimukseen valittiin viisi tiedontarvetta. Kullekkin tiedontarpeelle muodostettiin viisi kysymystä ja kukin näistä esitettiin järjestelmälle kymmenen kertaa. Yhteensä 250 vastausta arvioitiin manuaalisesti
  tuotettujen SQL-kyselyiden perusteella.
  Tiedontarpeen tulkinnan vaihtelu erotettiin saman tulkinnan
  vaihtoehtoisista SQL-toteutuksista.

  Suorituksista 173 tuotti oikean, 18 melkein oikean ja 19 väärän
  SQL-kyselyn. SQL-kysely jäi muodostamatta 40 suorituksessa. Selkeät
  kysymykset ja semanttisen näkymän valmiiden metriikoiden tukemat tehtävät
  tuottivat vakaita tulkintoja. Ongelmia ilmeni erityisesti implisiittisten
  taustaoletusten huomioimisessa, suodatuksissa ja duplikaattien
  käsittelyssä. Sanamuodon muutoksiin liittyi myös tulkintaeroja.

  Tulosten perusteella tulkinnan johdonmukaisuus ei takaa
  oikeellisuutta, eikä SQL-toteutusten vaihtelu itsessään osoita
  epäluotettavuutta. Järjestelmä tuotti tiedontarvetta vastaavia
  kyselyitä valituissa analyysitehtävissä, mutta oikeellisuus ei
  säilynyt kaikissa kysymysmuotoiluissa ja suorituksissa. Havainnot
  puoltavat kyselyiden oikeellisuuden varmistamista tutkitussa
  sovellusympäristössä. Tulokset ovat tapauskohtaisia eivätkä
  muodosta kattavaa arviota järjestelmän yleisestä suorituskyvystä.

  \vspace{1em}
  \noindent\textbf{Avainsanat:} kielimallit, Text-to-SQL, Cortex Analyst, Transformer, semanttinen näkymä, luotettavuus, autoregressiivinen mallinnus
\end{abstract}
